% Generated by scripts/assemble.py; corrections belong in source/corrections.json.
\examhead{29}{中国石油大学(华东) 2025 年数学分析试卷}{33}
% p034-b01; source PDF page 34
\blockanchor{p034-b01}
\question{1}求极限  $\lim_{x\to0}\cot x\left(\frac{1}{\sin x}-\frac{1}{x}\right)$ .\par
% p034-b02; source PDF page 34
\blockanchor{p034-b02}
\question{2}求不定积分 $\int x\arctan x\mathrm{d}x.$\par
% p034-b03; source PDF page 34
\blockanchor{p034-b03}
\question{3}计算第一类曲面积分 $\iint_{\Sigma} \frac{1}{z} \mathrm{~d} S$, 其中积分曲面 $\Sigma$ 是上半球面 $z = \sqrt{4 - x^2 - y^2}$ 被平面 $z = 1$ 所截下的顶部部分.\par
% p034-b04; source PDF page 34
\blockanchor{p034-b04}
\question{4}求幂级数 $\sum_{n=0}^{\infty} (-1)^n \frac{x^{2n+1}}{2n+1}$ 的收敛域及和函数 $S(x)$, 并求 $\sum_{n=0}^{\infty} (-1)^n \frac{1}{2n+1}$ 的和.\par
% p034-b05; source PDF page 34
\blockanchor{p034-b05}
\question{5}证明 $\{a_{n}\}$ 收敛的充要条件是级数 $\sum_{n=1}^{\infty}(a_{n+1} - a_{n})$ 收敛.\par
% p034-b06; source PDF page 34
\blockanchor{p034-b06}
\question{6}已知数列  $\{a_{n}\}$  非负有界, 证明:  $\lim_{n\to\infty}\sqrt[n]{a_{1}^{n}+a_{2}^{n}+\cdots+a_{n}^{n}}=\sup_{n\geq1}\{a_{n}\}$ .\par
% p034-b07; source PDF page 34
\blockanchor{p034-b07}
\question{7}证明: $\oint_{L} \frac{x \, \mathrm{d}y - y \, \mathrm{d}x}{3x^{2} + 4y^{2}} = \frac{\pi}{\sqrt{3}}$, 其中积分曲线 $L$: 椭圆 $2x^{2} + 3y^{2} = 1$, 方向为逆时针.\par
% p034-b08; source PDF page 34
\blockanchor{p034-b08}
\question{8}若 $f(x)$ 在 $[a, b]$ 内有有限导数，且 $f'(a) < 0$，$f'(b) > 0$，则至少存在 $c \in (a, b)$，使得 $f'(c) = 0$.\par
% p034-b09; source PDF page 34
\blockanchor{p034-b09}
\question{9}计算 $I = \iiint_{V} (xy^{2} + z) \, \mathrm{d}V$，其中 $V$ 是由曲面 $z = x^{2} + y^{2}$ 与平面 $z = 4$ 围成的区域.\par
% p034-b10; source PDF page 34
\blockanchor{p034-b10}
\question{10}证明: 若函数 $f(x)$ 在 $[a, b]$ 上可积, 函数 $g(x)$ 在 $[a, b]$ 上除一点 $x_0$ 外 $f(x) = g(x)$, 则 $g(x)$ 在 $[a, b]$ 上可积, 且 $\int_{a}^{b} f(x) \mathrm{d}x = \int_{a}^{b} g(x) \mathrm{d}x$.\par
% p034-b11; source PDF page 34
\blockanchor{p034-b11}
\question{11}分析函数 $f(x, y) = \begin{cases} \frac{x^4 + y^4}{x^2 + y^2}, & (x, y) \neq (0, 0); \\ 0, & (x, y) = (0, 0). \end{cases}$ 在 $(0, 0)$ 处的可微性.\par
% p034-b12; source PDF page 34
\blockanchor{p034-b12}
\question{12}设函数 $f(x)$ 在 [1,3] 上连续, 在 (1,3) 内可导, 且满足 $f(1) = \int_{2}^{3} x f(x) \, \mathrm{d}x$, 证明: 至少存在 $\eta \in (1,3)$, 使得 $f(\eta) + \eta f'(\eta) = 0$.\par
% p034-b13; source PDF page 34
\blockanchor{p034-b13}
\question{13}证明含参量积分 $F(y) = \int_{1}^{+\infty} \sin(xy)e^{-x^2} \, \mathrm{d}x$ 在 $y \in [1,2]$ 上可导.\par
% p034-b14; source PDF page 34
\par\Needspace{6\baselineskip}
\blockanchor{p034-b14}
\question{14}计算曲面积分\par
\nopagebreak[4]
% p034-b15; source PDF page 34
\blockanchor{p034-b15}
\displayformula{\iint_ {S} \frac {x ^ {2} \mathrm{d} y \mathrm{d} z + y ^ {2} \mathrm{d} z \mathrm{d} x + z ^ {2} \mathrm{d} x \mathrm{d} y}{\sqrt {x ^ {2} + y ^ {2} - z + 4}}.}
% p034-b16; source PDF page 34
\blockanchor{p034-b16}
其中 $S$ 是曲面 $x^{2} + y^{2} = z$ 介于平面 $z = 0$ 及 $z = 1$ 之间的部分, 方向取上侧.\par
% p034-b17; source PDF page 34
\blockanchor{p034-b17}
\question{15}用闭区间套定理证明致密性定理, 即设数列 $\{x_{n}\}$ 有界, 则必有收敛子列.\par
